\subsection{ROOTS}

The roots of G relative to T are the non-zero weights of the adjoint representation of G. The set of roots of G relative to T is denoted by R(G, T), or simply by R if there is no risk of confusion. By Prop. 6, the map

\[
\delta : \mathrm{X(T)} \to \mathfrak {t} _ {\mathbf {C}} ^ {*}
\]

\((t_{\mathbf{C}}^{*} \text{ denotes the dual of the complex vector space } t_{\mathbf{C}}) \text{ maps } R(G, T) \text{ bijectively onto the set } R(\mathfrak{g}_{\mathbf{C}}, t_{\mathbf{C}}) \text{ of roots of the split reductive algebra } (\mathfrak{g}_{\mathbf{C}}, t_{\mathbf{C}}) \text{ (Chap. VIII, §2, no. 2, Remark 4). If we put}\)

\[
\mathfrak {g} ^ {\alpha} = (\mathfrak {g} _ {\mathbf {C}}) _ {\alpha} (\mathrm{T}) = (\mathfrak {g} _ {\mathbf {C}}) _ {\delta (\alpha)} (\mathfrak {t} _ {\mathbf {C}}),\tag{11}
\]

for all \(\alpha\in R\) , then each \(\mathfrak{g}^{\alpha}\) is of dimension 1 over C (loc. cit., Th. 1) and

\[
\mathfrak {g} _ {\mathbf {C}} = t _ {\mathbf {C}} \oplus \bigoplus_ {\alpha \in R} \mathfrak {g} ^ {\alpha}.\tag{12}
\]

For each \(\alpha \in \mathbb{R}\) , denote by \(V(\alpha)\) the 2-dimensional subspace of \(\mathfrak{g}\) such that \(V(\alpha)_{(\mathbf{C})} = \mathfrak{g}^{\alpha} + \mathfrak{g}^{-\alpha}\) ; the non-zero isotypical components of \(\mathfrak{g}\) for the adjoint representation of T are \(\mathfrak{t}\) and the \(V(\alpha)\) . Further, let K be the quadratic form associated to the Killing form of \(\mathfrak{g}\) ; it is negative (§1, no. 3, Prop. 1) and its restriction \(\mathrm{K}(\alpha)\) to \(\mathrm{V}(\alpha)\) is negative and separating. For each element t of T, Ad t leaves \(\mathrm{K}(\alpha)\) stable, and hence gives a morphism of Lie groups

\[
\iota_ {\alpha}: \mathrm{T} \to \mathbf {S O} (\mathrm{K} (\alpha)).
\]

There exists a unique isomorphism \(\rho_{\alpha}:\mathbf{U}\to \mathbf{SO}(\mathrm{K}(\alpha))\) such that \(\iota_{\alpha} = \rho_{\alpha}\circ \alpha\) . Indeed, let \(X\) be a non-zero element of \(\mathfrak{g}^{\alpha}\) , and let \(Y\) be the image of \(X\) under the conjugation of \(\mathfrak{g}_{\mathbf{C}}\) relative to \(\mathfrak{g}\) ; then \(Y\in \mathfrak{g}^{-\alpha}\) , and we obtain a basis \((U,V)\) of \(\mathrm{V}(\alpha)\) by putting \(U = X + Y,V = i(X - Y)\) ; the matrix of the endomorphism of \(\mathrm{V}(\alpha)\) induced by \(\operatorname {Ad}t,t\in \mathrm{T}\) , with respect to the basis \((U,V)\) is

\[
\left( \begin{array}{c c} \mathcal {R} (t ^ {\alpha}) & - \mathcal {I} (t ^ {\alpha}) \\ \mathcal {I} (t ^ {\alpha}) & \mathcal {R} (t ^ {\alpha}) \end{array} \right),
\]

hence the assertion.

PROPOSITION 8. Let \(\mathrm{Q}(\mathrm{R})\) be the subgroup of \(\mathrm{X}(\mathrm{T})\) generated by the roots of \(\mathrm{G}\) .

\begin{enumerate}
\def\labelenumi{\alph{enumi})}
\item
  The centre \(\mathrm{C}(\mathrm{G})\) of G is a closed subgroup of T, equal to the intersection of the kernels of the roots. The canonical map \(\mathrm{X}(\mathrm{T}/\mathrm{C}(\mathrm{G})) \to \mathrm{X}(\mathrm{T})\) is injective with image \(\mathrm{Q}(\mathrm{R})\) .
\item
  The compact group C(G) is isomorphic to the dual of the discrete group X(T)/Q(R) (Spectral Theories, Chap. II, §1, no. 1, Def. 2).
\item
  C(G) reduces to the identity element if and only if Q(R) is equal to X(T).
\end{enumerate}

By §2, no. 2, Cor. 2 of Th. 2, C(G) is contained in T. Since this is the kernel of the adjoint representation, it is the intersection of the kernels of the roots, in other words the orthogonal complement of the subgroup Q(R) of X(T). Thus, the proposition follows from Spectral Theories, Chap. II, §1, no. 7, Th. 4 and no. 5, Th. 2.

PROPOSITION 9. Every automorphism of the Lie group G that induces the identity on T is of the form Int t, with \(t \in T\) .

Assume first of all that \(\mathrm{C}(\mathrm{G})\) reduces to the identity element, in other words that \(\mathrm{X}(\mathrm{T}) = \mathrm{Q}(\mathrm{R})\) (Prop. 8). Let f be an automorphism of G inducing the identity on T, and \(\varphi = \mathrm{L}(f)_{(\mathbf{C})}\) ; then \(\varphi\) is an automorphism of \(g_{C}\) inducing the identity on \(t_{C}\) . By Chap. VIII, §5, no. 2, Prop. 2, there exists a unique homomorphism \(\theta : \mathrm{Q}(\mathrm{R}) \to \mathbf{C}^{*}\) such that \(\varphi\) induces on each \(\mathfrak{g}^{\alpha}\) the homothety with ratio \(\theta(\alpha)\) . Since \(\varphi\) leaves stable the real form g of \(g_{C}\) , it commutes with the conjugation \(\sigma\) of \(g_{C}\) with respect to g; but \(\sigma(\mathfrak{g}^{\alpha}) = \mathfrak{g}^{-\alpha}\) , so \(\theta(-\alpha) = \overline{\theta(\alpha)}\) for all \(\alpha \in R\) . This implies that \(\theta(\alpha)\overline{\theta(\alpha)} = \theta(\alpha)\theta(-\alpha) = 1\) . It follows that \(\theta\) takes values in U, and hence corresponds by duality to an element t of T such that \((\operatorname{Ad} t)_{(\mathbf{C})} = \varphi\) , so \(\operatorname{Int} t = f\) .

In the general case, the preceding applies to the group G/C(G), whose centre reduces to the identity element, and to its maximal torus T/C(G).

It follows that, if \(f\) is an automorphism of \(G\) inducing the identity on \(T\) , there exists an element \(t\) of \(T\) such that \(f\) and \(\operatorname{Int} t\) induce by passage to the quotient the same automorphism of \(G / C(G)\) . But, since the canonical morphism \(D(G) \to G / C(G)\) is a finite covering (§1, no. 4, Cor. 1 of Prop. 4), \(f\) and \(\operatorname{Int} t\) induce the same automorphism of \(D(G)\) , hence of \(D(G) \times C(G)\) , and hence also of \(G\) (loc. cit.).

COROLLARY. Let u be an automorphism of G and H the closed subgroup of G consisting of the fixed points of u. Then, the automorphism u is inner if and only if \(H_{0}\) is of maximal rank.

If u is equal to Int g, with \(g \in G\) , the subgroup \(H_{0} = \mathrm{Z}(g)_{0}\) is of maximal rank ( \(\S2\) , no. 2, Cor. 3). Conversely, if H contains a maximal torus S, the automorphism u is of the form Int s with \(s \in S\) (Prop. 9).
% semantic-fragments: legacy-input-seam
\relax{}
